\section{Reference theorems and proof conventions}
\label{sec:foundations}

The probability model, its numerical parameters and the reporting scale are
assumptions or definitions. The conclusions below are mathematical consequences
of those choices. We use the cited standard foundational theorems; the proofs
here establish their particular forms needed by this report. In particular,
the cited Radon--Nikodym, spectral and characteristic-function uniqueness
theorems are not claims of new results. All random variables in a variance or
covariance calculation have finite second moments. Conditional-expectation
equalities hold almost surely unless a specified continuous version is being
evaluated at an individual history. Vector conditional expectations are taken
coordinate by coordinate.

\begin{theorem}[Moment identities and zero variance]
\label{thm:moments}
For integrable random variables, expectation is linear. If \(X_1,\ldots,X_n\)
have finite second moments, covariance is bilinear and
\[
  \Var\!\left(\sum_{j=1}^{n}a_jX_j\right)
  =\sum_{i=1}^{n}\sum_{j=1}^{n}a_i a_j\Cov(X_i,X_j).
\]
For a random vector \(X\), a constant matrix \(A\), and a constant vector
\(b\),
\[
  \E[AX+b]=A\E[X]+b,\qquad
  \Cov(AX+b)=A\Cov(X)A^{\T}.
\]
In particular, \(\operatorname{SD}(u+vX)=|v|\operatorname{SD}(X)\).
Moreover, \(\Var(X)=0\) if and only if \(X=\E[X]\) almost surely.
\end{theorem}

\begin{proof}
Linearity is the integral-linearity theorem
\cite[Theorem~1.6.1]{durrett2019}. Products of centered square-integrable
variables are integrable by H\"older's inequality
\cite[Theorem~1.6.3]{durrett2019}. Expanding
\(\E[(X-\E[X])(Y-\E[Y])]\), distributing finite sums and applying
linearity gives covariance bilinearity and the displayed scalar and matrix
identities. The SD identity follows by taking the nonnegative square root of
\(\Var(u+vX)=v^2\Var(X)\); see also the covariance summary in
\cite[Section~4.2]{bertsekas2008}. If the variance is zero, Chebyshev's
inequality gives
\[
  P\{|X-\E[X]|\geq\varepsilon\}
  \leq\frac{\Var(X)}{\varepsilon^2}=0
  \qquad(\varepsilon>0)
\]
\cite[Theorem~1.6.4]{durrett2019}. Taking the union of these zero-probability
events for \(\varepsilon=1,1/2,1/3,\ldots\) proves the assertion. Its
converse follows directly from the variance definition.
\end{proof}

\begin{theorem}[Conditional expectation and total variance]
\label{thm:conditional}
Let \(\mathcal G\) be a sub-\(\sigma\)-algebra of the probability space and
let \(X\) be integrable. There exists a \(\mathcal G\)-measurable integrable
variable \(Z=\E[X\given\mathcal G]\), unique up to almost-sure equality,
such that
\[
  \E[\one_A Z]=\E[\one_A X]\qquad\text{for every }A\in\mathcal G.
\]
Conditional expectation is linear, fixes integrable \(\mathcal G\)-measurable
variables, and equals \(\E[X]\) when \(X\) is independent of
\(\mathcal G\). If \(\mathcal G\subseteq\mathcal H\), its tower property is
\[
  \E[\E[X\given\mathcal H]\given\mathcal G]
  =\E[X\given\mathcal G].
\]
If \(H\) is \(\mathcal G\)-measurable and \(X,HX\) are integrable, the
pull-out identity is \(\E[HX\given\mathcal G]=H\E[X\given\mathcal G]\).
For square-integrable \(X\),
\[
  \Var(X)=\Var(\E[X\given\mathcal G])
      +\E[\Var(X\given\mathcal G)],
\]
where
\(\Var(X\given\mathcal G)
 =\E[(X-\E[X\given\mathcal G])^2\given\mathcal G]\).
\end{theorem}

\begin{proof}
Existence follows by applying the Radon--Nikodym theorem separately to the
finite positive measures \(A\mapsto\E[\one_A X^+]\) and
\(A\mapsto\E[\one_A X^-]\) on \(\mathcal G\). Both are absolutely
continuous with respect to the probability measure restricted to
\(\mathcal G\). Subtracting their derivatives gives the defining integral
identity; uniqueness gives almost-sure uniqueness. This is the
construction in \cite[Section~4.1 and Theorem~A.4.8]{durrett2019}.
Linearity and the known-variable property follow by substituting candidates
in that identity. Independence gives
\(\E[\one_A X]=P(A)\E[X]\), proving the independent-variable case.
For the tower property, integrate over \(A\in\mathcal G\subseteq\mathcal H\)
twice and invoke uniqueness. These properties are respectively recorded in
\cite[Examples~4.1.3--4.1.4 and Theorems~4.1.9(a), 4.1.13]{durrett2019}.
The pull-out property follows first for simple \(H\), then by approximation
and truncation under the stated integrability assumptions
\cite[Theorem~4.1.14]{durrett2019}.

For total variance put \(Z=\E[X\given\mathcal G]\) and \(R=X-Z\).
The conditional-expectation projection theorem gives \(Z\in L^2\) and
\(\E[RH]=0\) for every \(\mathcal G\)-measurable \(H\in L^2\)
\cite[Theorem~4.1.15]{durrett2019}. Apply it to \(H=Z-\E[X]\) and expand
the square in \(X-\E[X]=R+(Z-\E[X])\). The cross term vanishes, while
the tower property gives
\(\E[R^2]=\E[\E[R^2\given\mathcal G]]\). This proves total variance,
also stated in \cite[Section~4.3]{bertsekas2008} and posed in
\cite[Exercise~4.1.7]{durrett2019}.
\end{proof}

\begin{theorem}[Spectral criterion and Gaussian covariance construction]
\label{thm:spectral}
A real symmetric matrix \(S\) is positive semidefinite if and only if all
its eigenvalues are nonnegative, and positive definite if and only if all
are positive. Every positive-semidefinite matrix is a covariance matrix of
a possibly singular Gaussian vector. Every covariance matrix is positive
semidefinite.
\end{theorem}

\begin{proof}
The real spectral theorem gives \(S=U\Lambda U^{\T}\) with \(U\)
orthogonal and \(\Lambda\) real diagonal
\cite[Theorem~7.29]{axler2024}. Hence
\(t^{\T}St=\sum_j\lambda_j(U^{\T}t)_j^2\), proving both eigenvalue
criteria. For nonnegative eigenvalues let
\(S^{1/2}=U\operatorname{diag}(\sqrt{\lambda_j})U^{\T}\), the positive
square root \cite[Theorem~7.38]{axler2024}. If \(G\) has independent
standard-normal coordinates, the vector \(S^{1/2}G\) has covariance \(S\)
by Theorem~\ref{thm:moments}; this construction includes zero eigenvalues
\cite[Theorems~3.4.4--3.4.5, author-hosted draft]{gallager2013}.
Conversely, \(t^{\T}\Cov(X)t=\Var(t^{\T}X)\geq0\) for every \(t\).
\end{proof}

\begin{theorem}[Gaussian transformations and block independence]
\label{thm:gaussian}
An affine transformation of a jointly Gaussian vector is jointly Gaussian,
including when its covariance is singular. Its mean and covariance determine
its law. If two jointly Gaussian blocks have zero cross-covariance, they are
independent, even if either block is singular or constant.
\end{theorem}

\begin{proof}
Write a Gaussian vector as \(X=\mu+AG\), where \(G\) has independent
standard-normal coordinates; \(A\) need not be invertible. Then
\(BX+b=(B\mu+b)+BAG\), proving closure
\cite[Definition~3.3.1 and Theorem~3.3.1, author-hosted draft]{gallager2013}.
The one-dimensional Gaussian characteristic function and independence of
the coordinates of \(G\) give
\[
  \phi_X(t)=\E[e^{\mathrm i t^{\T}X}]
    =\exp\!\left(\mathrm i t^{\T}\mu-\tfrac12t^{\T}\Sigma t\right),
  \qquad\Sigma=AA^{\T},
\]
as in \cite[Theorem~3.3.2 and equations~(3.14), (3.18), author-hosted draft]{gallager2013}.
The uniqueness consequence of multivariate characteristic-function inversion
\cite[Theorem~3.10.4]{durrett2019} shows that \(\mu,\Sigma\) determine
the distribution. For blocks \((U,V)\) with \(\Cov(U,V)=0\), the displayed
formula factors as \(\phi_{(U,V)}(s,t)=\phi_U(s)\phi_V(t)\).
By the product-measure construction and Fubini's theorem
\cite[Theorems~1.7.1--1.7.2]{durrett2019}, this is the characteristic function
of the product of the two marginal laws. Uniqueness therefore proves independence.
This proof uses no density
and no covariance inverse, so it includes singular blocks.
\end{proof}

\begin{theorem}[Gaussian regression with a nonsingular conditioning block]
\label{thm:regression}
Let \((U,V)\) be jointly Gaussian with means \(\mu_U,\mu_V\) and covariance
blocks \(A=\Cov(U)\), \(B=\Cov(V)\), \(C=\Cov(U,V)\). Suppose only
that \(B\) is positive definite. A conditional distribution of \(U\) given
\(V=v\) is the possibly singular normal law with
\[
  m(v)=\mu_U+CB^{-1}(v-\mu_V),\qquad
  K=A-CB^{-1}C^{\T}.
\]
In particular its conditional mean is \(m(v)\) and its conditional covariance
is \(K\). The full joint covariance need not be nonsingular.
\end{theorem}

\begin{proof}
Set \(W=U-\mu_U-CB^{-1}(V-\mu_V)\). The moment identities yield
\[
  \E[W]=0,\qquad \Cov(W,V)=C-CB^{-1}B=0,
  \qquad \Cov(W)=A-CB^{-1}C^{\T}.
\]
Theorem~\ref{thm:gaussian} shows that \(W\) is Gaussian and independent of
\(V\). Thus \(U=m(V)+W\). Independence implies, for every bounded Borel
function \(h\), that a version of \(\E[h(U)\given V]\) is
\(v\mapsto\int h(m(v)+w)\,P_W(dw)\). Independence and Fubini's theorem
\cite[Theorem~1.7.2]{durrett2019} verify its defining conditional integral
identity. This conditional probability kernel has the stated normal law,
mean and covariance. The argument
also works when \(K\) has zero eigenvalues.
The familiar full-rank result is given in
\cite[Theorems~3.5.1--3.5.2 and equations~(3.44)--(3.45), author-hosted draft]{gallager2013}.
Those cited statements assume a nonsingular full joint covariance; the
residual proof here supplies the version needed when only the conditioning
block is nonsingular.
\end{proof}

\begin{lemma}[Unique continuous conditional-mean version]
\label{lem:continuous-version}
If \(V\) is a Gaussian vector with positive-definite covariance, two
continuous functions representing the same conditional mean given \(V\)
agree at every point. In particular, the affine mean function in
Theorem~\ref{thm:regression} is its unique continuous version.
\end{lemma}

\begin{proof}
The nonsingular Gaussian density theorem
\cite[Theorem~3.4.6 and equation~(3.22), author-hosted draft]{gallager2013}
gives, in dimension \(q\),
\[
  f_V(v)=\frac{\exp[-\tfrac12(v-\mu_V)^{\T}B^{-1}(v-\mu_V)]}
                 {(2\pi)^{q/2}\sqrt{\det B}}>0
  \qquad\text{for every }v\in\mathbb R^q.
\]
Theorem~\ref{thm:conditional} makes any two conditional-mean versions equal
\(P_V\)-almost everywhere. If continuous versions differed at one point,
one coordinate of their difference would remain nonzero on some open ball.
Such a ball has positive probability under the displayed density, contradicting
almost-sure equality. This justifies evaluating the specified continuous
version at a point history; it does not turn a probability-zero event into
a probability ratio.
\end{proof}

\begin{theorem}[Mathematical induction]
\label{thm:induction}
For an assertion \(P(n)\) on integers \(n\geq n_0\), the base assertion
\(P(n_0)\) and the implications \(P(n)\Rightarrow P(n+1)\) for every
\(n\geq n_0\) establish \(P(n)\) for all such integers
\cite[Section~10.1]{hammack2018}.
\end{theorem}

\begin{proof}
If there were a counterexample, the well-ordering principle would supply
the least one. It could not be \(n_0\), by the base assertion. Its predecessor
would satisfy \(P\), and the induction step would then rule out this least
counterexample. This is a contradiction.
\end{proof}
